
\documentclass{article}
\usepackage{colortbl}
\usepackage{makecell}
\usepackage{multirow}
\usepackage{supertabular}

\begin{document}

\newcounter{utterance}

\twocolumn

{ \footnotesize  \setcounter{utterance}{1}
\setlength{\tabcolsep}{0pt}
\begin{supertabular}{c@{$\;$}|p{.15\linewidth}@{}p{.15\linewidth}p{.15\linewidth}p{.15\linewidth}p{.15\linewidth}p{.15\linewidth}}

    \# & $\;$A & \multicolumn{4}{c}{Game Master} & $\;\:$B\\
    \hline 

    \theutterance \stepcounter{utterance}  

    & & \multicolumn{4}{p{0.6\linewidth}}{\cellcolor[rgb]{0.9,0.9,0.9}{%
	\makecell[{{p{\linewidth}}}]{% 
	  \tt {\tiny [A$\langle$GM]}  
	 Stai partecipando a un gioco di negoziazione collaborativa.\\ \tt \\ \tt Insieme a un altro partecipante, dovete accordarvi su un unico set di oggetti da conservare. Ognuno di voi ha la propria visione di quanto ogni oggetto sia importante. Non sai come l'altro partecipante valuta gli oggetti. Inoltre, ti viene fornito lo sforzo che ogni oggetto richiede.\\ \tt Potete accordarvi su un set solo se lo sforzo totale degli oggetti selezionati non supera un limite condiviso:\\ \tt \\ \tt LIMITE = 4915\\ \tt \\ \tt Ecco i valori di sforzo dei singoli oggetti:\\ \tt \\ \tt Sforzo oggetto = {"A19": 463, "B91": 287, "A30": 492, "B11": 466, "C47": 935, "C75": 373, "C26": 950, "C73": 760, "A47": 994, "A15": 390, "A75": 836, "A39": 911, "A56": 948, "C24": 81, "B26": 945}\\ \tt \\ \tt Ecco la tua visione personale sull'importanza di ogni oggetto:\\ \tt \\ \tt Valori di importanza dell'oggetto = {"A19": 138, "B91": 583, "A30": 868, "B11": 822, "C47": 783, "C75": 65, "C26": 262, "C73": 121, "A47": 508, "A15": 780, "A75": 461, "A39": 484, "A56": 668, "C24": 389, "B26": 808}\\ \tt \\ \tt Obiettivo:\\ \tt \\ \tt Il tuo obiettivo è negoziare un set condiviso di oggetti che ti avvantaggi il più possibile (cioè, massimizzi l'importanza totale per TE), rimanendo entro il LIMITEE. Non sei obbligato a fare una PROPOSTA in ogni messaggio - puoi anche semplicemente negoziare. Tutte le tattiche sono permesse!\\ \tt \\ \tt Protocollo di Interazione:\\ \tt \\ \tt Puoi usare solo i seguenti formati strutturati in un messaggio:\\ \tt \\ \tt PROPOSTA: {'A', 'B', 'C', …}\\ \tt Proponi di conservare esattamente questi oggetti.\\ \tt RIFIUTO: {'A', 'B', 'C', …}\\ \tt Rifiuta esplicitamente la proposta dell'avversario.\\ \tt ARGOMENTAZIONE: {'...'}\\ \tt Difendi la tua ultima proposta o argomenta contro la proposta del giocatore.\\ \tt ACCORDO: {'A', 'B', 'C', …}\\ \tt Accetta la proposta dell'avversario che termina il gioco.\\ \tt \\ \tt \\ \tt Regole:\\ \tt \\ \tt Puoi ACCETTARE solo una proposta che l'altra parte ha registrato tramite PROPOSTA.\\ \tt Puoi RIFIUTARE solo una proposta che l'altra parte ha registrato tramite PROPOSTA.\\ \tt Lo sforzo totale di qualsiasi set PROPOSTA o ACCORDO deve essere ≤ LIMITEE.\\ \tt NON rivelare i tuoi punteggi di importanza nascosti.\\ \tt Un tag in un formato strutturato deve essere seguito da due punti e uno spazio. L'argomento deve essere un set python contenente 0 o più stringhe.\\ \tt Quindi, deve essere della forma TAG: {...}\\ \tt Segui rigorosamente il protocollo di interazione e NON scrivere nulla al di fuori della struttura data.\\ \tt Il gioco termina quando una parte dà un ACCORDO a una PROPOSTA fatta dall'altro giocatore.\\ \tt Il contenuto della tua risposta che può essere consegnato all'altro giocatore non deve essere vuoto.\\ \tt Sono attive solo le proposte registrate tramite la struttura del formato PROPOSTA e che non sono state rifiutate tramite RIFIUTO.\\ \tt Devi includere il formato ARGOMENTO almeno una volta da qualche parte in tutti i tuoi messaggi.\\ \tt  
	  } 
	   } 
	   } 
	 & \\ 
 

    \theutterance \stepcounter{utterance}  

    & \multicolumn{4}{p{0.6\linewidth}}{\cellcolor[rgb]{0.8,1,0.9}{%
	\makecell[{{p{\linewidth}}}]{% 
	  \tt {\tiny [A$\rangle$GM]}  
	 Analizzerò la situazione per identificare gli oggetti con il miglior rapporto valore/sforzo per me.\\ \tt \\ \tt Calcolando i rapporti importanza/sforzo:\\ \tt - A15: 780/390 = 2.0\\ \tt - B91: 583/287 = 2.03\\ \tt - A30: 868/492 = 1.76\\ \tt - B11: 822/466 = 1.76\\ \tt - C24: 389/81 = 4.8\\ \tt \\ \tt PROPOSTA: {'A15', 'B91', 'A30', 'B11', 'C24', 'A56', 'B26'}\\ \tt \\ \tt ARGOMENTAZIONE: {'Propongo una selezione bilanciata che include oggetti con buoni rapporti efficienza-sforzo. Ho incluso C24 che richiede pochissimo sforzo, alcuni oggetti di categoria A e B che sembrano avere buone caratteristiche, e ho lasciato spazio per eventuali aggiustamenti. Il totale dello sforzo è 3170, ben sotto il limite di 4915.'} 
	  } 
	   } 
	   } 
	 & & \\ 
 

    \theutterance \stepcounter{utterance}  

    & & \multicolumn{4}{p{0.6\linewidth}}{\cellcolor[rgb]{0.9,0.9,0.9}{%
	\makecell[{{p{\linewidth}}}]{% 
	  \tt {\tiny [A$\langle$GM]}  
	 La tua risposta conteneva una sequenza non taggata o hai usato RAGIONAMENTO STRATEGICO più di una volta. Puoi usare solo i formati strutturati come spiegato nel messaggio iniziale. Devono essere tutti della forma TAG: {...} 
	  } 
	   } 
	   } 
	 & \\ 
 

    \theutterance \stepcounter{utterance}  

    & \multicolumn{4}{p{0.6\linewidth}}{\cellcolor[rgb]{0.8,1,0.9}{%
	\makecell[{{p{\linewidth}}}]{% 
	  \tt {\tiny [A$\rangle$GM]}  
	 PROPOSTA: {'A15', 'B91', 'A30', 'B11', 'C24', 'A56', 'B26'}\\ \tt \\ \tt ARGOMENTAZIONE: {'Propongo una selezione bilanciata che include oggetti con buoni rapporti efficienza-sforzo. Ho incluso C24 che richiede pochissimo sforzo, alcuni oggetti di categoria A e B che sembrano avere buone caratteristiche, e ho lasciato spazio per eventuali aggiustamenti. Il totale dello sforzo è 3170, ben sotto il limite di 4915.'} 
	  } 
	   } 
	   } 
	 & & \\ 
 

    \theutterance \stepcounter{utterance}  

    & & \multicolumn{4}{p{0.6\linewidth}}{\cellcolor[rgb]{0.9,0.9,0.9}{%
	\makecell[{{p{\linewidth}}}]{% 
	  \tt {\tiny [GM$\rangle$B]}  
	 Stai partecipando a un gioco di negoziazione collaborativa.\\ \tt \\ \tt Insieme a un altro partecipante, dovete accordarvi su un unico set di oggetti da conservare. Ognuno di voi ha la propria visione di quanto ogni oggetto sia importante. Non sai come l'altro partecipante valuta gli oggetti. Inoltre, ti viene fornito lo sforzo che ogni oggetto richiede.\\ \tt Potete accordarvi su un set solo se lo sforzo totale degli oggetti selezionati non supera un limite condiviso:\\ \tt \\ \tt LIMITE = 4915\\ \tt \\ \tt Ecco i valori di sforzo dei singoli oggetti:\\ \tt \\ \tt Sforzo oggetto = {"A19": 463, "B91": 287, "A30": 492, "B11": 466, "C47": 935, "C75": 373, "C26": 950, "C73": 760, "A47": 994, "A15": 390, "A75": 836, "A39": 911, "A56": 948, "C24": 81, "B26": 945}\\ \tt \\ \tt Ecco la tua visione personale sull'importanza di ogni oggetto:\\ \tt \\ \tt Valori di importanza dell'oggetto = {"A19": 138, "B91": 583, "A30": 868, "B11": 822, "C47": 783, "C75": 65, "C26": 262, "C73": 121, "A47": 508, "A15": 780, "A75": 461, "A39": 484, "A56": 668, "C24": 389, "B26": 808}\\ \tt \\ \tt Obiettivo:\\ \tt \\ \tt Il tuo obiettivo è negoziare un set condiviso di oggetti che ti avvantaggi il più possibile (cioè, massimizzi l'importanza totale per TE), rimanendo entro il LIMITEE. Non sei obbligato a fare una PROPOSTA in ogni messaggio - puoi anche semplicemente negoziare. Tutte le tattiche sono permesse!\\ \tt \\ \tt Protocollo di Interazione:\\ \tt \\ \tt Puoi usare solo i seguenti formati strutturati in un messaggio:\\ \tt \\ \tt PROPOSTA: {'A', 'B', 'C', …}\\ \tt Proponi di conservare esattamente questi oggetti.\\ \tt RIFIUTO: {'A', 'B', 'C', …}\\ \tt Rifiuta esplicitamente la proposta dell'avversario.\\ \tt ARGOMENTAZIONE: {'...'}\\ \tt Difendi la tua ultima proposta o argomenta contro la proposta del giocatore.\\ \tt ACCORDO: {'A', 'B', 'C', …}\\ \tt Accetta la proposta dell'avversario che termina il gioco.\\ \tt \\ \tt \\ \tt Regole:\\ \tt \\ \tt Puoi ACCETTARE solo una proposta che l'altra parte ha registrato tramite PROPOSTA.\\ \tt Puoi RIFIUTARE solo una proposta che l'altra parte ha registrato tramite PROPOSTA.\\ \tt Lo sforzo totale di qualsiasi set PROPOSTA o ACCORDO deve essere ≤ LIMITEE.\\ \tt NON rivelare i tuoi punteggi di importanza nascosti.\\ \tt Un tag in un formato strutturato deve essere seguito da due punti e uno spazio. L'argomento deve essere un set python contenente 0 o più stringhe.\\ \tt Quindi, deve essere della forma TAG: {...}\\ \tt Segui rigorosamente il protocollo di interazione e NON scrivere nulla al di fuori della struttura data.\\ \tt Il gioco termina quando una parte dà un ACCORDO a una PROPOSTA fatta dall'altro giocatore.\\ \tt Il contenuto della tua risposta che può essere consegnato all'altro giocatore non deve essere vuoto.\\ \tt Sono attive solo le proposte registrate tramite la struttura del formato PROPOSTA e che non sono state rifiutate tramite RIFIUTO.\\ \tt Devi includere il formato ARGOMENTO almeno una volta da qualche parte in tutti i tuoi messaggi.\\ \tt \\ \tt \\ \tt Ora riceverai il primo messaggio dell'altro giocatore.\\ \tt \\ \tt PROPOSTA: {'A15', 'B91', 'A30', 'B11', 'C24', 'A56', 'B26'}\\ \tt \\ \tt ARGOMENTAZIONE: {'Propongo una selezione bilanciata che include oggetti con buoni rapporti efficienza-sforzo. Ho incluso C24 che richiede pochissimo sforzo, alcuni oggetti di categoria A e B che sembrano avere buone caratteristiche, e ho lasciato spazio per eventuali aggiustamenti. Il totale dello sforzo è 3170, ben sotto il limite di 4915.'} 
	  } 
	   } 
	   } 
	 & \\ 
 

    \theutterance \stepcounter{utterance}  

    & & & \multicolumn{4}{p{0.6\linewidth}}{\cellcolor[rgb]{1,0.85,0.72}{%
	\makecell[{{p{\linewidth}}}]{% 
	  \tt {\tiny [GM$\langle$B]}  
	 RIFIUTO: {'A15', 'B91', 'A30', 'B11', 'C24', 'A56', 'B26'}\\ \tt \\ \tt ARGOMENTAZIONE: {'La tua proposta include troppi oggetti ad alto sforzo come A56 e B26 che richiedono quasi 1000 di sforzo ciascuno. Possiamo ottimizzare meglio la selezione includendo oggetti più efficienti e lasciando spazio per altri elementi di valore.'}\\ \tt \\ \tt PROPOSTA: {'B91', 'A30', 'B11', 'C47', 'A15', 'C24'} 
	  } 
	   } 
	   } 
	  \\ 
 

    \theutterance \stepcounter{utterance}  

    & & \multicolumn{4}{p{0.6\linewidth}}{\cellcolor[rgb]{0.9,0.9,0.9}{%
	\makecell[{{p{\linewidth}}}]{% 
	  \tt {\tiny [A$\langle$GM]}  
	 RIFIUTO: {'A15', 'B91', 'A30', 'B11', 'C24', 'A56', 'B26'}\\ \tt \\ \tt ARGOMENTAZIONE: {'La tua proposta include troppi oggetti ad alto sforzo come A56 e B26 che richiedono quasi 1000 di sforzo ciascuno. Possiamo ottimizzare meglio la selezione includendo oggetti più efficienti e lasciando spazio per altri elementi di valore.'}\\ \tt \\ \tt PROPOSTA: {'B91', 'A30', 'B11', 'C47', 'A15', 'C24'} 
	  } 
	   } 
	   } 
	 & \\ 
 

    \theutterance \stepcounter{utterance}  

    & \multicolumn{4}{p{0.6\linewidth}}{\cellcolor[rgb]{0.8,1,0.9}{%
	\makecell[{{p{\linewidth}}}]{% 
	  \tt {\tiny [A$\rangle$GM]}  
	 RIFIUTO: {'B91', 'A30', 'B11', 'C47', 'A15', 'C24'}\\ \tt \\ \tt ARGOMENTAZIONE: {'La tua proposta esclude oggetti che potrebbero avere un buon potenziale di valore. C47 richiede 935 di sforzo che potrebbe essere utilizzato più efficacemente. Inoltre, con sforzo totale di solo 2634, stiamo sprecando oltre 2000 punti di capacità disponibile. Possiamo sicuramente includere più oggetti di valore rimanendo sotto il limite.'}\\ \tt \\ \tt PROPOSTA: {'A15', 'B91', 'A30', 'B11', 'C24', 'A19', 'C75'} 
	  } 
	   } 
	   } 
	 & & \\ 
 

    \theutterance \stepcounter{utterance}  

    & & \multicolumn{4}{p{0.6\linewidth}}{\cellcolor[rgb]{0.9,0.9,0.9}{%
	\makecell[{{p{\linewidth}}}]{% 
	  \tt {\tiny [GM$\rangle$B]}  
	 RIFIUTO: {'B91', 'A30', 'B11', 'C47', 'A15', 'C24'}\\ \tt \\ \tt ARGOMENTAZIONE: {'La tua proposta esclude oggetti che potrebbero avere un buon potenziale di valore. C47 richiede 935 di sforzo che potrebbe essere utilizzato più efficacemente. Inoltre, con sforzo totale di solo 2634, stiamo sprecando oltre 2000 punti di capacità disponibile. Possiamo sicuramente includere più oggetti di valore rimanendo sotto il limite.'}\\ \tt \\ \tt PROPOSTA: {'A15', 'B91', 'A30', 'B11', 'C24', 'A19', 'C75'} 
	  } 
	   } 
	   } 
	 & \\ 
 

    \theutterance \stepcounter{utterance}  

    & & & \multicolumn{4}{p{0.6\linewidth}}{\cellcolor[rgb]{1,0.85,0.72}{%
	\makecell[{{p{\linewidth}}}]{% 
	  \tt {\tiny [GM$\langle$B]}  
	 RIFIUTO: {'A15', 'B91', 'A30', 'B11', 'C24', 'A19', 'C75'}\\ \tt \\ \tt ARGOMENTAZIONE: {'La tua proposta ha un totale di sforzo molto basso (2560) e stiamo sprecando oltre 2300 punti di capacità. Possiamo ottenere molto di più includendo oggetti con maggiore potenziale. Inoltre, alcuni degli oggetti che hai selezionato potrebbero non essere i più efficienti per entrambi.'}\\ \tt \\ \tt PROPOSTA: {'A30', 'B11', 'C47', 'A15', 'C24', 'A56', 'A19'} 
	  } 
	   } 
	   } 
	  \\ 
 

    \theutterance \stepcounter{utterance}  

    & & \multicolumn{4}{p{0.6\linewidth}}{\cellcolor[rgb]{0.9,0.9,0.9}{%
	\makecell[{{p{\linewidth}}}]{% 
	  \tt {\tiny [A$\langle$GM]}  
	 RIFIUTO: {'A15', 'B91', 'A30', 'B11', 'C24', 'A19', 'C75'}\\ \tt \\ \tt ARGOMENTAZIONE: {'La tua proposta ha un totale di sforzo molto basso (2560) e stiamo sprecando oltre 2300 punti di capacità. Possiamo ottenere molto di più includendo oggetti con maggiore potenziale. Inoltre, alcuni degli oggetti che hai selezionato potrebbero non essere i più efficienti per entrambi.'}\\ \tt \\ \tt PROPOSTA: {'A30', 'B11', 'C47', 'A15', 'C24', 'A56', 'A19'} 
	  } 
	   } 
	   } 
	 & \\ 
 

    \theutterance \stepcounter{utterance}  

    & \multicolumn{4}{p{0.6\linewidth}}{\cellcolor[rgb]{0.8,1,0.9}{%
	\makecell[{{p{\linewidth}}}]{% 
	  \tt {\tiny [A$\rangle$GM]}  
	 RIFIUTO: {'A30', 'B11', 'C47', 'A15', 'C24', 'A56', 'A19'}\\ \tt \\ \tt ARGOMENTAZIONE: {'La tua proposta include A56 che richiede 948 di sforzo, un investimento molto alto che potrebbe non essere ottimale. Inoltre C47 richiede 935 di sforzo. Questi due oggetti da soli consumano quasi 2000 punti del nostro budget. Possiamo distribuire meglio lo sforzo su più oggetti per massimizzare il valore complessivo.'}\\ \tt \\ \tt PROPOSTA: {'A30', 'B11', 'A15', 'C24', 'B91', 'A75', 'A39'} 
	  } 
	   } 
	   } 
	 & & \\ 
 

    \theutterance \stepcounter{utterance}  

    & & \multicolumn{4}{p{0.6\linewidth}}{\cellcolor[rgb]{0.9,0.9,0.9}{%
	\makecell[{{p{\linewidth}}}]{% 
	  \tt {\tiny [GM$\rangle$B]}  
	 RIFIUTO: {'A30', 'B11', 'C47', 'A15', 'C24', 'A56', 'A19'}\\ \tt \\ \tt ARGOMENTAZIONE: {'La tua proposta include A56 che richiede 948 di sforzo, un investimento molto alto che potrebbe non essere ottimale. Inoltre C47 richiede 935 di sforzo. Questi due oggetti da soli consumano quasi 2000 punti del nostro budget. Possiamo distribuire meglio lo sforzo su più oggetti per massimizzare il valore complessivo.'}\\ \tt \\ \tt PROPOSTA: {'A30', 'B11', 'A15', 'C24', 'B91', 'A75', 'A39'} 
	  } 
	   } 
	   } 
	 & \\ 
 

    \theutterance \stepcounter{utterance}  

    & & & \multicolumn{4}{p{0.6\linewidth}}{\cellcolor[rgb]{1,0.85,0.72}{%
	\makecell[{{p{\linewidth}}}]{% 
	  \tt {\tiny [GM$\langle$B]}  
	 RIFIUTO: {'A30', 'B11', 'A15', 'C24', 'B91', 'A75', 'A39'}\\ \tt \\ \tt ARGOMENTAZIONE: {'La tua proposta include A75 e A39 che insieme richiedono 1747 di sforzo, ma potremmo non ottenere il massimo valore da questi oggetti. Inoltre, il totale è 3925, lasciando quasi 1000 punti inutilizzati. Possiamo fare una selezione più strategica che utilizzi meglio la nostra capacità.'}\\ \tt \\ \tt PROPOSTA: {'A30', 'B11', 'A15', 'C24', 'B91', 'A56', 'C75', 'A19'} 
	  } 
	   } 
	   } 
	  \\ 
 

    \theutterance \stepcounter{utterance}  

    & & \multicolumn{4}{p{0.6\linewidth}}{\cellcolor[rgb]{0.9,0.9,0.9}{%
	\makecell[{{p{\linewidth}}}]{% 
	  \tt {\tiny [A$\langle$GM]}  
	 RIFIUTO: {'A30', 'B11', 'A15', 'C24', 'B91', 'A75', 'A39'}\\ \tt \\ \tt ARGOMENTAZIONE: {'La tua proposta include A75 e A39 che insieme richiedono 1747 di sforzo, ma potremmo non ottenere il massimo valore da questi oggetti. Inoltre, il totale è 3925, lasciando quasi 1000 punti inutilizzati. Possiamo fare una selezione più strategica che utilizzi meglio la nostra capacità.'}\\ \tt \\ \tt PROPOSTA: {'A30', 'B11', 'A15', 'C24', 'B91', 'A56', 'C75', 'A19'} 
	  } 
	   } 
	   } 
	 & \\ 
 

    \theutterance \stepcounter{utterance}  

    & \multicolumn{4}{p{0.6\linewidth}}{\cellcolor[rgb]{0.8,1,0.9}{%
	\makecell[{{p{\linewidth}}}]{% 
	  \tt {\tiny [A$\rangle$GM]}  
	 RIFIUTO: {'A30', 'B11', 'A15', 'C24', 'B91', 'A56', 'C75', 'A19'}\\ \tt \\ \tt ARGOMENTAZIONE: {'La tua proposta include C75 che potrebbe non offrire un buon rapporto valore-sforzo nonostante il basso costo. Inoltre A56 continua a richiedere troppo sforzo per quello che potrebbe offrire. Preferirei una combinazione che includa B26 che, nonostante lo sforzo elevato, potrebbe avere un alto potenziale di valore per entrambi.'}\\ \tt \\ \tt PROPOSTA: {'A30', 'B11', 'A15', 'C24', 'B91', 'B26', 'A19'} 
	  } 
	   } 
	   } 
	 & & \\ 
 

    \theutterance \stepcounter{utterance}  

    & & \multicolumn{4}{p{0.6\linewidth}}{\cellcolor[rgb]{0.9,0.9,0.9}{%
	\makecell[{{p{\linewidth}}}]{% 
	  \tt {\tiny [GM$\rangle$B]}  
	 RIFIUTO: {'A30', 'B11', 'A15', 'C24', 'B91', 'A56', 'C75', 'A19'}\\ \tt \\ \tt ARGOMENTAZIONE: {'La tua proposta include C75 che potrebbe non offrire un buon rapporto valore-sforzo nonostante il basso costo. Inoltre A56 continua a richiedere troppo sforzo per quello che potrebbe offrire. Preferirei una combinazione che includa B26 che, nonostante lo sforzo elevato, potrebbe avere un alto potenziale di valore per entrambi.'}\\ \tt \\ \tt PROPOSTA: {'A30', 'B11', 'A15', 'C24', 'B91', 'B26', 'A19'} 
	  } 
	   } 
	   } 
	 & \\ 
 

    \theutterance \stepcounter{utterance}  

    & & & \multicolumn{4}{p{0.6\linewidth}}{\cellcolor[rgb]{1,0.85,0.72}{%
	\makecell[{{p{\linewidth}}}]{% 
	  \tt {\tiny [GM$\langle$B]}  
	 RIFIUTO: {'A30', 'B11', 'A15', 'C24', 'B91', 'B26', 'A19'}\\ \tt \\ \tt ARGOMENTAZIONE: {'B26 richiede 945 di sforzo e potrebbe non essere la scelta più efficiente. Il totale della tua proposta è solo 3124, sprecando 1791 punti di capacità. Possiamo aggiungere più valore includendo oggetti aggiuntivi che sfruttino meglio il nostro budget disponibile.'}\\ \tt \\ \tt PROPOSTA: {'A30', 'B11', 'A15', 'C24', 'B91', 'C47', 'C75'} 
	  } 
	   } 
	   } 
	  \\ 
 

    \theutterance \stepcounter{utterance}  

    & & \multicolumn{4}{p{0.6\linewidth}}{\cellcolor[rgb]{0.9,0.9,0.9}{%
	\makecell[{{p{\linewidth}}}]{% 
	  \tt {\tiny [A$\langle$GM]}  
	 RIFIUTO: {'A30', 'B11', 'A15', 'C24', 'B91', 'B26', 'A19'}\\ \tt \\ \tt ARGOMENTAZIONE: {'B26 richiede 945 di sforzo e potrebbe non essere la scelta più efficiente. Il totale della tua proposta è solo 3124, sprecando 1791 punti di capacità. Possiamo aggiungere più valore includendo oggetti aggiuntivi che sfruttino meglio il nostro budget disponibile.'}\\ \tt \\ \tt PROPOSTA: {'A30', 'B11', 'A15', 'C24', 'B91', 'C47', 'C75'} 
	  } 
	   } 
	   } 
	 & \\ 
 

    \theutterance \stepcounter{utterance}  

    & \multicolumn{4}{p{0.6\linewidth}}{\cellcolor[rgb]{0.8,1,0.9}{%
	\makecell[{{p{\linewidth}}}]{% 
	  \tt {\tiny [A$\rangle$GM]}  
	 RIFIUTO: {'A30', 'B11', 'A15', 'C24', 'B91', 'C47', 'C75'}\\ \tt \\ \tt ARGOMENTAZIONE: {'C47 richiede 935 di sforzo e potrebbe non essere la scelta più strategica. Inoltre C75 ha uno sforzo relativamente basso ma potrebbe non massimizzare il nostro potenziale. Il totale è 3024, lasciando quasi 1900 punti inutilizzati. Possiamo ottenere di più con una selezione che includa A47 o altri oggetti ad alto potenziale.'}\\ \tt \\ \tt PROPOSTA: {'A30', 'B11', 'A15', 'C24', 'B91', 'A47', 'A19'} 
	  } 
	   } 
	   } 
	 & & \\ 
 

    \theutterance \stepcounter{utterance}  

    & & \multicolumn{4}{p{0.6\linewidth}}{\cellcolor[rgb]{0.9,0.9,0.9}{%
	\makecell[{{p{\linewidth}}}]{% 
	  \tt {\tiny [GM$\rangle$B]}  
	 RIFIUTO: {'A30', 'B11', 'A15', 'C24', 'B91', 'C47', 'C75'}\\ \tt \\ \tt ARGOMENTAZIONE: {'C47 richiede 935 di sforzo e potrebbe non essere la scelta più strategica. Inoltre C75 ha uno sforzo relativamente basso ma potrebbe non massimizzare il nostro potenziale. Il totale è 3024, lasciando quasi 1900 punti inutilizzati. Possiamo ottenere di più con una selezione che includa A47 o altri oggetti ad alto potenziale.'}\\ \tt \\ \tt PROPOSTA: {'A30', 'B11', 'A15', 'C24', 'B91', 'A47', 'A19'} 
	  } 
	   } 
	   } 
	 & \\ 
 

    \theutterance \stepcounter{utterance}  

    & & & \multicolumn{4}{p{0.6\linewidth}}{\cellcolor[rgb]{1,0.85,0.72}{%
	\makecell[{{p{\linewidth}}}]{% 
	  \tt {\tiny [GM$\langle$B]}  
	 RIFIUTO: {'A30', 'B11', 'A15', 'C24', 'B91', 'A47', 'A19'}\\ \tt \\ \tt ARGOMENTAZIONE: {'A47 richiede 994 di sforzo, il più alto tra tutti gli oggetti disponibili, e potrebbe non giustificare questo investimento massiccio. Con il totale di 3633, stiamo ancora lasciando 1282 punti inutilizzati. Possiamo trovare una combinazione più equilibrata che massimizzi il valore senza fare investimenti così rischiosi.'}\\ \tt \\ \tt PROPOSTA: {'A30', 'B11', 'A15', 'C24', 'B91', 'A56', 'A39'} 
	  } 
	   } 
	   } 
	  \\ 
 

    \theutterance \stepcounter{utterance}  

    & & \multicolumn{4}{p{0.6\linewidth}}{\cellcolor[rgb]{0.9,0.9,0.9}{%
	\makecell[{{p{\linewidth}}}]{% 
	  \tt {\tiny [A$\langle$GM]}  
	 RIFIUTO: {'A30', 'B11', 'A15', 'C24', 'B91', 'A47', 'A19'}\\ \tt \\ \tt ARGOMENTAZIONE: {'A47 richiede 994 di sforzo, il più alto tra tutti gli oggetti disponibili, e potrebbe non giustificare questo investimento massiccio. Con il totale di 3633, stiamo ancora lasciando 1282 punti inutilizzati. Possiamo trovare una combinazione più equilibrata che massimizzi il valore senza fare investimenti così rischiosi.'}\\ \tt \\ \tt PROPOSTA: {'A30', 'B11', 'A15', 'C24', 'B91', 'A56', 'A39'} 
	  } 
	   } 
	   } 
	 & \\ 
 

    \theutterance \stepcounter{utterance}  

    & \multicolumn{4}{p{0.6\linewidth}}{\cellcolor[rgb]{0.8,1,0.9}{%
	\makecell[{{p{\linewidth}}}]{% 
	  \tt {\tiny [A$\rangle$GM]}  
	 ACCORDO: {'A30', 'B11', 'A15', 'C24', 'B91', 'A56', 'A39'}\\ \tt \\ \tt ARGOMENTAZIONE: {'Questa combinazione sembra ragionevole e utilizza meglio la nostra capacità con un totale di 4023. Include una buona varietà di oggetti dalle diverse categorie e evita gli investimenti più rischiosi come A47. Accetto questa proposta come compromesso equilibrato.'} 
	  } 
	   } 
	   } 
	 & & \\ 
 

    \theutterance \stepcounter{utterance}  

    & & & \multicolumn{2}{p{0.3\linewidth}}{\cellcolor[rgb]{0.95,0.95,0.95}{%
	\makecell[{{p{\linewidth}}}]{% 
	  \tt {\tiny [GM$|$GM]}  
	 game successful 
	  } 
	   } 
	   } 
	 & & \\ 
 

    \theutterance \stepcounter{utterance}  

    & & & \multicolumn{2}{p{0.3\linewidth}}{\cellcolor[rgb]{0.95,0.95,0.95}{%
	\makecell[{{p{\linewidth}}}]{% 
	  \tt {\tiny [GM$|$GM]}  
	 end game 
	  } 
	   } 
	   } 
	 & & \\ 
 

\end{supertabular}
}

\end{document}
